% (User input window)
% 2026-02-15 20:12:08（Asia/Singapore）
% According to an internal note dated 2026-02-14.

\subsection{Determinant Potentials of the Fiber Spectrum: Multiplicity Operator and Reflector Projection as Readouts at the Same Temperature}\label{subsec:op-algebra-fiber-determinant-potentials}
Proposition \ref{prop:fold-fiber-multiplicity-trace-moments} has shown: the folding fiber multiplicities \(\{d_m(x)\}_{x\in X_m}\) can be lifted to a diagonal spectral operator \(\mathcal{F}_m\) on the stable layer, whose trace moments \(\Tr(\mathcal{F}_m^q)\) are precisely the collision moments \(S_q(m)\).
This subsection records a complementary spectral functor to the trace moments: the determinant potential generated by the same spectral data \(\{d_m(x)\}\), which in regime terms compresses "microscopic volume, boundary volume, and fiber spectrum" into readouts of the same partition function at different temperatures.

\begin{proposition}[Fiber Spectrum Partition Function: Closed Form and Two-Endpoint Readout of \texorpdfstring{$Z_m(t)=\det(I+t\mathcal{F}_m)$}{Zm(t)=det(I+tF)}]\label{prop:op-algebra-fiber-spectrum-determinant-partition}
Under the notation of Definition \ref{def:fold-fiber-multiplicity-operator}, for any \(t\ge 0\), define the determinant potential
\[
Z_m(t):=\det(I+t\mathcal{F}_m).
\]
Then the complete closed form holds
\[
\boxed{\;
Z_m(t)=\prod_{x\in X_m}\bigl(1+t\,d_m(x)\bigr)\;}
\]
and its low-temperature/high-temperature limits satisfy
\[
\boxed{\;
\lim_{t\downarrow 0}\frac{Z_m(t)-1}{t}=\Tr(\mathcal{F}_m)=\sum_{x\in X_m}d_m(x)=|\Omega_m|=2^m\;}
\]
as well as
\[
\boxed{\;
\lim_{t\to\infty}\frac{\dd}{\dd(\log t)}\log Z_m(t)
=\lim_{t\to\infty}\sum_{x\in X_m}\frac{t\,d_m(x)}{1+t\,d_m(x)}
=|X_m|\;}
\]
thus
\[
\boxed{\;
\lim_{t\to\infty}\Bigl(\log Z_m(t)-|X_m|\log t\Bigr)=\sum_{x\in X_m}\log d_m(x)\;}
\]
yielding a strict interpolation readout that simultaneously encapsulates the boundary volume \(|X_m|\) and the fiber spectrum \(\{d_m(x)\}\) into a single determinant potential.
\end{proposition}
\begin{proof}
By Proposition \ref{prop:fold-fiber-multiplicity-trace-moments}, \(\mathcal{F}_m\) is diagonal in the standard basis with eigenvalues \(\{d_m(x)\}_{x\in X_m}\), hence
\(
Z_m(t)=\prod_x(1+t d_m(x))
\).
First-order expansion as \(t\downarrow 0\) yields
\(
Z_m(t)=1+t\Tr(\mathcal{F}_m)+O(t^2)
\);
while \(\sum_x d_m(x)=|\Omega_m|\) is the fiber cardinality summation identity.
\par
For the derivative, directly from
\(
\log Z_m(t)=\sum_x \log(1+t d_m(x))
\)
we obtain
\(
\frac{\dd}{\dd(\log t)}\log Z_m(t)=\sum_x \frac{t d_m(x)}{1+t d_m(x)}
\to |X_m|
\) (termwise limit).
The last identity follows from \(\log(1+t d)=\log t+\log(d+1/t)\) and letting \(t\to\infty\). QED.
\end{proof}

\begin{corollary}[Determinant Potential of the Fiber Reflector: \texorpdfstring{$\det(I+tK_m)=(1+t)^{|X_m|}$}{det(I+tK)=(1+t)^|X|}]\label{cor:op-algebra-fiber-reflector-determinant-potential}
Viewing the fiber reflector \(K_m\) as a linear operator on \(\CC^{\Omega_m}\) (Theorem \ref{thm:pom-fiber-reflector} and Corollary \ref{cor:pom-Km-spectrum}), for any \(t>-1\) we have
\[
\boxed{\;
\det(I+tK_m)=(1+t)^{|X_m|},\qquad
\frac{1}{m}\log\det(I+tK_m)=\frac{|X_m|}{m}\log(1+t)\;}
\]
thus the exponential growth rate of the boundary volume \(|X_m|\) can be losslessly read out from the normalized determinant potential at any fixed \(t\neq 0\).
\end{corollary}
\begin{proof}
By Corollary \ref{cor:pom-Km-spectrum}, the spectrum of \(K_m\) consists only of \(\{0,1\}\) with eigenvalue \(1\) having multiplicity \(|X_m|\).
Thus the eigenvalues of \(I+tK_m\) are \(1+t\) (multiplicity \(|X_m|\)) and \(1\) (multiplicity \(2^m-|X_m|\)), taking the determinant yields the conclusion. QED.
\end{proof}

